% Chapter 10 worked example. Included by ch10_resources.tex.
\section{Worked Example: Patient-Facing Chatbot and Digital Inclusion}
\label{sec:comp-ch10-worked-example}

\subsection{Purpose and Status}
This worked example demonstrates how the methods in Chapter~10 can be
translated into a reviewable institutional decision. All numerical inputs are
synthetic unless a source is identified explicitly. They are intended to show the
logic of the analysis, not to report empirical findings or establish a universal
threshold. Local use requires replacement of the assumptions with current,
versioned evidence and approval through the relevant clinical, managerial,
economic, legal, technical, and governance processes.

A hospital proposes a multilingual patient-facing chatbot for appointment guidance, preparation instructions, and symptom navigation. The intervention could improve access but may also create factual, privacy, cybersecurity, accessibility, workload, and trust risks. 

\subsection{Decision Problem and Comparators}
\textbf{Decision problem.} Should the chatbot be piloted, restricted to administrative tasks, redesigned with human escalation, or rejected in favour of non-AI navigation improvements?

The comparator set is deliberately broader than the existing pathway. This prevents
a new technology from receiving a lower evidentiary threshold than staffing,
workflow, shared-service, conventional digital, or no-change alternatives.

\begin{longtable}{@{}p{0.08\textwidth}p{0.84\textwidth}@{}}
\caption{Comparators for the Chapter 10 worked example}
\label{tab:comp-ch10-comparators}\\
\toprule
\textbf{Option} & \textbf{Comparator or strategy} \\
\midrule
\endfirsthead
\toprule
\textbf{Option} & \textbf{Comparator or strategy} \\
\midrule
\endhead
\bottomrule
\endlastfoot
1 & Improved static web and printed information \\
2 & Call-centre redesign and multilingual staffing \\
3 & Rules-based conversational interface \\
4 & Restricted AI chatbot with human escalation \\
\end{longtable}

\subsection{Model and Decision Structure}
The analytical structure should follow the decision rather than the available
software. The team should map the causal pathway from intervention input to
information, human decision, action, outcome, resource consequence, and review.
The structure should also identify capacity constraints, uncertainty, subgroup
variation, implementation requirements, and events that would change the
recommendation.

A compact quantitative relationship used in this example is:
\begin{equation}
\text{No compensatory risk equation is recommended; mandatory safeguards are assessed as gates.}
\label{eq:comp-ch10-key-relationship}
\end{equation}
The equation is an organizing aid rather than a substitute for disaggregated evidence,
qualitative findings, or mandatory safeguards. Its variables and interpretation must
be defined in the local protocol.

\subsection{Synthetic Parameters and Evidence Sources}
Table~\ref{tab:comp-ch10-parameters} provides the base-case inputs. A
production analysis should add parameter distributions, correlations, structural
alternatives, data dates, system versions, and evidence-confidence ratings.

\begin{longtable}{@{}p{0.32\textwidth}p{0.22\textwidth}p{0.38\textwidth}@{}}
\caption{Synthetic parameters for the Chapter 10 worked example}
\label{tab:comp-ch10-parameters}\\
\toprule
\textbf{Parameter} & \textbf{Base value} & \textbf{Source or interpretation} \\
\midrule
\endfirsthead
\toprule
\textbf{Parameter} & \textbf{Base value} & \textbf{Source or interpretation} \\
\midrule
\endhead
\bottomrule
\endlastfoot
Languages proposed & 8 & Service plan \\
Share of users requiring assisted digital access & 18\% & Patient survey \\
Queries classified as clinically consequential & 24\% & Retrospective sample \\
Human escalation target & Under 5 minutes & Safety requirement \\
Retention of interaction logs & 30 days & Draft policy \\
Expected annual interactions & 420,000 & Activity forecast \\
\end{longtable}

\subsection{Step-by-Step Analysis}
The sequence below is intended to make the analysis reproducible and to ensure
that evidence, economics, governance, and implementation develop together.

\begin{longtable}{@{}p{0.07\textwidth}p{0.55\textwidth}p{0.30\textwidth}@{}}
\caption{Analytical sequence}
\label{tab:comp-ch10-analysis-steps}\\
\toprule
\textbf{Step} & \textbf{Required work} & \textbf{Decision record} \\
\midrule
\endfirsthead
\toprule
\textbf{Step} & \textbf{Required work} & \textbf{Decision record} \\
\midrule
\endhead
\bottomrule
\endlastfoot
1 & Define allowed, prohibited, and escalation-required tasks. & Evidence and responsible owner to be recorded \\
2 & Map data flows, third parties, retention, logging, and model-service dependencies. & Evidence and responsible owner to be recorded \\
3 & Assess factuality, information integrity, privacy, cybersecurity, language, accessibility, and digital inclusion. & Evidence and responsible owner to be recorded \\
4 & Estimate the effect on call-centre demand, repeat contacts, abandonment, and downstream services. & Evidence and responsible owner to be recorded \\
5 & Assign accountability, incident communication, contestability, and stopping authority. & Evidence and responsible owner to be recorded \\
6 & Pilot only where essential non-digital routes remain available and monitoring can detect differential use and harm. & Evidence and responsible owner to be recorded \\
\end{longtable}

\subsection{Illustrative Outputs}
The outputs in Table~\ref{tab:comp-ch10-outputs} show how technical,
clinical, operational, economic, equity, and governance findings should be kept
visible rather than compressed into a single favourable or unfavourable score.

\begin{longtable}{@{}p{0.30\textwidth}p{0.24\textwidth}p{0.38\textwidth}@{}}
\caption{Illustrative model and decision outputs}
\label{tab:comp-ch10-outputs}\\
\toprule
\textbf{Output} & \textbf{Result} & \textbf{Interpretation} \\
\midrule
\endfirsthead
\toprule
\textbf{Output} & \textbf{Result} & \textbf{Interpretation} \\
\midrule
\endhead
\bottomrule
\endlastfoot
Suitable initial scope & Administrative guidance and verified preparation instructions & Symptom advice remains restricted \\
Primary equity risk & Exclusion of users with language, disability, literacy, connectivity, or trust barriers & Requires human alternatives \\
Primary governance risk & Unclear responsibility for harmful or incomplete advice & Requires explicit escalation and logs \\
Expected call reduction & Uncertain & Must account for repeat and escalated contacts \\
Preferred decision & Restricted pilot with independent content verification & Not unrestricted patient-facing deployment \\
\end{longtable}

\subsection{Sensitivity, Uncertainty, and Subgroups}
At minimum, the completed analysis should vary the parameters most likely to
change the decision, test at least one credible alternative model structure, and
report subgroup results separately where performance, access, response, burden,
or opportunity cost may differ. Deterministic analysis should identify thresholds;
probabilistic analysis should be used where credible parameter distributions and
correlations are available. Scenario analysis is preferable when structural or
future uncertainty cannot be represented defensibly by one probability model.

The record should distinguish uncertainty that can be reduced through evidence,
uncertainty that can be managed through safeguards, and uncertainty that remains
too consequential to accept. Missing evidence should not be represented as an
average or neutral value without justification.

\subsection{Interpretation and Recommendation}
Pilot a restricted, clearly disclosed service for low-consequence tasks, preserve immediate human alternatives, and prohibit unsupported diagnosis or treatment advice. Expansion should depend on effective access and outcome evidence rather than interaction volume.

The recommendation is conditional rather than permanent. It applies only to the
specified population, setting, intervention configuration, evidence cut-off, and
version. The following conditions should be incorporated into the approval,
contract, implementation, monitoring, or renewal record:

\begin{itemize}
 \item all content domains have accountable clinical or operational owners.
 \item language and accessibility testing includes affected users.
 \item privacy, retention, and third-party data use are transparent and minimized.
 \item escalation reliability and abandonment are monitored by subgroup.
 \item a serious factual, privacy, security, or equity failure triggers immediate restriction and review.
\end{itemize}

\subsection{Decision Statement}
\begin{quote}
For the specified decision problem and comparator set, the institution should
\emph{[proceed / proceed with conditions / redesign / pilot / restrict / defer /
replace / retire]}. The decision applies to \emph{[population, setting, version,
and evidence period]}, is owned by \emph{[accountable authority]}, and will be
reviewed on \emph{[date]} or earlier if \emph{[trigger events]} occur. The
evidence, assumptions, unresolved uncertainty, mandatory safeguards, monitoring
thresholds, and exit arrangements are recorded in the accompanying dossier.
\end{quote}
